\chapter{结论}
\label{cha:conclusion}

\section*{工作总结}

本文设计并实现了一个Python程序执行路径追踪和代码覆盖率分析工具——PathTracer。该工具具有以下特点：

\textbf{1. 技术实现}

\begin{itemize}
    \item 使用Python标准库\texttt{sys.settrace}实现运行时追踪，记录程序执行的每一行代码
    \item 使用\texttt{ast}模块进行静态分析，识别源代码中的所有控制流分支点
    \item 通过对比静态分支与动态执行路径，计算代码覆盖率
\end{itemize}

\textbf{2. 功能特性}

\begin{itemize}
    \item 支持识别顺序、选择（if/elif/else）、循环（for/while）、异常处理（try/except）等程序结构
    \item 记录函数调用树，包含参数和返回值信息
    \item 支持文本、JSON、HTML三种报告输出格式
    \item 提供命令行接口和Python API两种使用方式
\end{itemize}

\textbf{3. 实验验证}

通过完整的实验验证，工具成功：

\begin{itemize}
    \item 设计并执行了3个测试用例（TC1/TC2/TC3），测试套件累计覆盖率达72.7\%
    \item 对demo.py进行了完整测试，覆盖了if/for/while/except等多种控制流结构
    \item 与coverage.py进行了对比分析（使用--branch模式），明确了两者统计口径差异
    \item 生成了HTML可视化报告（见\texttt{docs/coverage\_demo\_full.html}）
\end{itemize}

\textbf{4. 工具定位}

本实验验证了PathTracer在Python分支覆盖与执行路径跟踪上的可行性，适用于以下场景：

\begin{itemize}
    \item \textbf{教学演示}：帮助理解程序控制流和执行路径
    \item \textbf{调试分析}：定位未覆盖的代码分支，辅助测试设计
    \item \textbf{代码审查}：记录函数调用关系和参数传递
\end{itemize}

\textbf{局限性}：本实验主要针对小型示例程序进行验证，尚未针对大型真实项目（如Django、Flask等框架）做系统评测。工具在大规模代码库上的性能和准确性有待进一步验证。

\section*{工作展望}

未来可以从以下方面改进工具：

\begin{enumerate}
    \item \textbf{提高覆盖率精度}：增加对elif分支、else分支的独立识别，区分true/false两个执行方向
    \item \textbf{支持更多语法}：扩展支持with语句、async/await、match-case等Python 3.10+新特性
    \item \textbf{集成开发环境}：开发IDE插件，实现实时覆盖率显示
    \item \textbf{增量分析}：支持增量覆盖率分析，只分析变更的代码
    \item \textbf{测试用例关联}：将覆盖率与具体测试用例关联，便于定位测试盲点
    \item \textbf{大规模验证}：在真实大型项目上进行系统评测，验证工具的泛化能力
\end{enumerate}
